% Generated draft: requires mathematical and source review.
\documentclass[11pt]{article}
\usepackage[margin=25mm]{geometry}
\usepackage{fontspec}
\setmainfont{Latin Modern Roman}
\usepackage{amsmath,amssymb,mathtools,unicode-math}
\setmathfont{Latin Modern Math}
\usepackage{enumitem,needspace,xurl,hyperref}
\setlength{\parindent}{0pt}
\setlength{\parskip}{5pt}
\setlength{\emergencystretch}{3em}
\newcommand{\N}{\mathbb{N}}
\newcommand{\Z}{\mathbb{Z}}
\newcommand{\Q}{\mathbb{Q}}
\newcommand{\R}{\mathbb{R}}
\newcommand{\C}{\mathbb{C}}
\newcommand{\F}{\mathbb{F}}
\newcommand{\A}{\mathbb{A}}
\newcommand{\PP}{\mathbb{P}}
\newcommand{\E}{\mathbb{E}}
\newcommand{\eps}{\varepsilon}
\newcommand{\dd}{\,\mathrm{d}}
\begin{document}
\section*{UnsolvedMath ID: 30005530}
\textbf{Source problem code:} catalogue entry
\subsection{Definitions and mathematical statement}
Let \(F\) be a holomorphic Siegel Hecke cusp form of degree two and weight \(k\) for \(\operatorname{Sp}_4(\Z)\), as in the supplied statement. Write
\[
\mathfrak H_2=\{Z\in M_2(\C):Z^{\mathsf T}=Z,\ \operatorname{Im}Z>0\}.
\]
Use the conventional Fourier-index lattice
\[
\Lambda=\left\{S=\begin{pmatrix}a&b/2\\b/2&c\end{pmatrix}:a,b,c\in\Z,\ a>0,\ 4ac-b^2>0\right\}.
\]
Thus \(S\) is positive definite and half-integral. Its associated integral binary quadratic form is \(ax^2+bxy+cy^2\), with discriminant
\[
D(S)=b^2-4ac=-4\det S.
\]
Let \(\Lambda_0\subseteq\Lambda\) consist of the indices for which \(D(S)\) is a negative fundamental discriminant, meaning the discriminant of an imaginary quadratic field. These definitions reconstruct the conventional meaning of the source's notation; the supplied excerpt does not explicitly define \(\Lambda\) or \(\Lambda_0\). Define the coefficients by
\[
F(Z)=\sum_{S\in\Lambda}a(F,S)\exp\!\bigl(2\pi i\operatorname{tr}(SZ)\bigr),\qquad Z\in\mathfrak H_2.
\]

The first requested assertion is that, for each fixed \(F\) as above, there exist \(\delta>0\) and \(C_F>0\), independent of \(S\), such that
\[
|a(F,S)|\le C_F(\det S)^{(k-1)/2-\delta}\qquad\text{for every }S\in\Lambda_0.
\]
The source calls \(\delta\) fixed but does not specify whether it must also be independent of \(F\) and \(k\). The displayed formulation therefore records the fixed-\(F\) assertion and leaves the stronger uniform interpretation unresolved. The source explicitly permits a conditional treatment using standard conjectures on \(L\)-functions, including the generalized Riemann hypothesis (GRH), and expected period identities. Any such dependencies must be stated.

The stronger conjectural target, attributed in the source to Reznikov--Saldana, is
\[
\forall\epsilon>0\ \exists C_{F,\epsilon}>0\ \forall S\in\Lambda_0:\qquad
|a(F,S)|\le C_{F,\epsilon}(\det S)^{k/2-3/4+\epsilon}.
\]
The source additionally proposes this stronger estimate for every \(S\in\Lambda\), provided that \(F\) is not a Saito--Kurokawa lift. That exclusion is supplied for the extension to \(\Lambda\); the supplied fundamental-discriminant formulation does not impose it.

For context, the supplied excerpt attributes to Kohnen the unconditional estimate
\[
|a(F,S)|\ll_{F,\epsilon}(\det S)^{k/2-13/36+\epsilon}\qquad(S\in\Lambda_0,\ \epsilon>0).
\]
Its description as the current best unconditional bound is a source claim, not an independently established literature assessment here. The excerpt also attributes a conditional estimate under GRH to Jaasaari--Lester--Saha. The apparent reading of its damaged typesetting is
\[
|a(F,S)|\ll_{F,\epsilon}
\frac{(\det S)^{k/2-1/2}}{(\log(\det S))^{1/8}}.
\]
This transcription requires checking. In particular, although the excerpt says that the estimate holds for all \(S\in\Lambda_0\), the displayed logarithmic denominator does not define a positive real expression when \(\det S\le1\); no small-determinant convention is supplied. It is therefore recorded only as an apparent large-determinant reading, not as a precise all-index theorem.

The excerpt motivates these estimates by an expected period identity following from the Gan--Gross--Prasad conjectures. It involves, for a class-group character \(\chi\), a sum of the form
\[
\sum_{\substack{[S]\in\operatorname{SL}_2(\Z)\backslash\Lambda_0\\D(S)=d}}
\chi([S])a(F,S)
\]
and the central value \(L(1/2,F\times\theta_\chi)\), where the source does not define the normalization of \(\theta_\chi\). The extracted formula contains a power resembling \(d^{k-1}\), but its placement of a square, its use of signed or absolute discriminant, and its proportionality factors are unclear. No precise period identity is asserted in this draft.
\subsection{Short English statement}
For a fixed holomorphic Siegel Hecke cusp form of degree two, seek a power saving in its Fourier coefficients at positive definite indices of fundamental discriminant. The stronger target is the supplied Reznikov--Saldana estimate, with a proposed extension to all positive definite indices after excluding Saito--Kurokawa lifts. Conditional arguments are explicitly permitted. Uniformity of the saving, damaged background formulas, and the supplied literature assessment require review.
\subsection{Sources}
\begin{itemize}
\item Supplied local UnsolvedMath registry record, permanent ID \(30005530\), title \emph{Ramanujan Bounds for Siegel Cusp-Form Coefficients}, exact source problem code \texttt{catalogue entry}, in the Oberwolfach Reports Open Problems collection. No individual proposer is supplied; the proposed year is 2023.
\item Supplied report citation: \emph{Group Actions and Harmonic Analysis in Number Theory} (2023), Oberwolfach Reports, DOI \texttt{10.4171/owr/2023/22}, pp.\ 1195--1262. \url{https://doi.org/10.4171/owr/2023/22}. The excerpt identifies the primary request as Question 3, the Kohnen estimate as Theorem 6, the conditional Jaasaari--Lester--Saha estimate as Theorem 7, and the Reznikov--Saldana target as Conjecture 2.
\item Supplied literature source labeled ``Blomer--Pohl theorem'': \url{https://arxiv.org/abs/1402.4635}.
\item Supplied literature source labeled ``published sup-norm paper'': \url{https://doi.org/10.1353/ajm.2016.0032}.
\item Literature assessment checked at the exact supplied date \texttt{2026-08-22}: the registry labels the problem \texttt{partially\_solved} and states that power-saving sup-norm bounds hold on each fixed compact set for degree-two Hecke--Maass forms, while a stronger global uniform bound remains unresolved. This is supplied dated curation, not independently verified. Its subject differs from the holomorphic Fourier-coefficient problem formulated here, so it does not establish partial resolution of this statement.
\item The registry labels the statement \texttt{corrected\_verified} and supplies the assertion that it was reconstructed and checked against the official Oberwolfach report. This is a supplied verification claim; no independent source consultation or mathematical verification is claimed here. The background explicitly requests independent verification of time-sensitive status and cited claims.
\end{itemize}
\end{document}
